% Generated draft: requires mathematical and source review.
\documentclass[11pt]{article}
\usepackage[margin=25mm]{geometry}
\usepackage{fontspec}
\setmainfont{Latin Modern Roman}
\usepackage{amsmath,amssymb,mathtools,unicode-math}
\setmathfont{Latin Modern Math}
\usepackage{enumitem,needspace,xurl,hyperref}
\setlength{\parindent}{0pt}
\setlength{\parskip}{5pt}
\setlength{\emergencystretch}{3em}
\newcommand{\N}{\mathbb{N}}
\newcommand{\Z}{\mathbb{Z}}
\newcommand{\Q}{\mathbb{Q}}
\newcommand{\R}{\mathbb{R}}
\newcommand{\C}{\mathbb{C}}
\newcommand{\F}{\mathbb{F}}
\newcommand{\A}{\mathbb{A}}
\newcommand{\PP}{\mathbb{P}}
\newcommand{\E}{\mathbb{E}}
\newcommand{\eps}{\varepsilon}
\newcommand{\dd}{\,\mathrm{d}}
\begin{document}
\section*{UnsolvedMath ID: 6000014}
\textbf{Source problem code:} catalogue entry
\subsection{Definitions and mathematical statement}
The proposed evolution of a plane curve is
\[
\frac{\partial x}{\partial u}=(1+kp)x'',
\]
where $u$ is time, $x''$ is the affine normal, $k$ is affine curvature, and $p$ is affine support function. The source asks for a time-global solution.

In a common nondegenerate equiaffine convention, affine arclength $s$ is chosen so that $\det(x_s,x_{ss})=1$; the affine normal is $x_{ss}$, and curvature is defined by $x_{sss}=-k x_s$. These conventions identify one possible meaning of the primes, but are not specified in the source. The affine support function additionally depends on the chosen origin and determinant/sign convention; the record supplies no formula for $p$. That sign affects the factor $1+kp$ and cannot be chosen silently.

A forward global initial-value formulation would specify a curve class $\mathcal C$ and require, for every $x_0\in\mathcal C$, an evolution $x(\cdot,u)$ defined for every $u\geq0$, satisfying the equation and $x(\cdot,0)=x_0$, in a prescribed regularity class and with the required nondegeneracy or convexity retained. The source gives neither $\mathcal C$ nor an initial condition, curve parameter domain, boundary/periodicity conditions, or the intended time interval. Mere existence of some eternal or stationary solution is different from this uniform initial-value assertion.

The time derivative must also be specified relative to a fixed material parameter or a time-dependent affine-arclength gauge; tangential reparametrization can change the written evolution. No global-existence theorem, preservation property, or normalization preventing degeneracy or extinction is supplied, so these remain formulation requirements rather than inferred facts.
\subsection{Short English statement}
Establish global existence for the affine curve flow $\partial_u x=(1+kp)x''$. Specify the initial curves, affine support convention, parametrization, and meaning of global regularity before asserting a precise theorem.
\subsection{Sources}
\begin{itemize}
\item Local UnsolvedMath registry, permanent ID 6000014, exact source code catalogue entry, ``Global Affine Curve Flow''; proposed by Takeshi Sasaki, 1998. Registry status partially solved, classification PARTIAL-PROGRESS.
\item Furuhata, Matsuzoe, and Urakawa, \emph{Open Problems in Affine Differential Geometry and Related Topics} (1998), item 4(c), scanned PDF pages 1--2, \url{https://www.jstage.jst.go.jp/article/iis/4/2/4_2_125/_pdf/-char/en}. Supplied manual OCR from page images accessed July 29, 2026. No PDF or image was accessed here.
\item The supplied research summary notes existence and long-time results for related affine shortening and central affine evolutions but says it did not verify an explicit global-existence theorem for this particular equation. No precise later reference is supplied, and related flows are not substituted for it.
\item Provenance says the historical article presents the problem as open and flags current status NEEDS\_REVIEW. The supplied rights note marks the J-STAGE article CC BY 4.0; attribution is retained.
\end{itemize}
\end{document}
